\documentclass[aps,prl,reprint,nofootinbib,superscriptaddress]{revtex4-2}
% Review Letter -- "Accessory = Node": a spectral-dynamical identity for integrable Landau-Zener.
% Memorializes the R8 theorem (v_* = E_*). Companion: type1_lz_working_paper.tex; proof artifacts
% experiments/r8_proof.py, paper/notes/r8_accessory_node_proof.md. Evidence: PROVEN (symbolic identity
% over Q(gamma,eps,a) via gauge slice + covariance; independently re-verified in exact arithmetic).
\usepackage{amsmath,amssymb,bm}
\newcommand{\eps}{\varepsilon}
\newcommand{\Scat}{\mathcal{S}}
\newcommand{\Pmid}{P_{m\to m}}
\newcommand{\Z}{\mathbb{Q}}
\newcommand{\diag}{\mathrm{diag}}
\newcommand{\Res}{\mathrm{Res}}
\newcommand{\Disc}{\mathrm{Disc}}

\begin{document}

\title{Pinning the accessory parameter: a spectral--dynamical identity\\
for the integrable multistate Landau--Zener problem}
\author{Co-theoretical-physicist working session}
\affiliation{\texttt{projects/type-1-lz}; branch \texttt{claude/comathematician-agent-skills-9HPSJ}}
\date{\today}

\begin{abstract}
The transition amplitudes of a multistate Landau--Zener (MLZ) model are the Stokes data of an
irregular ordinary differential equation; for the integrable Type-1 family at $N{=}3$ that ODE
carries a single finite \emph{apparent} singularity whose location is the \emph{accessory parameter}
--- the one datum that, in the general theory of such connection problems, is a free transcendental
modulus and is precisely what forbids a closed-form answer. We prove that for every generic Type-1
$N{=}3$ model this accessory parameter is \emph{not} free: it coincides exactly, as a rational
function of the couplings, with the energy of the model's guaranteed exact level crossing,
$v_*\equiv E_*$ in $\Z(\gamma,\eps,a)$. The proof has two faces: an elimination identity
($\Res$-based) and a structural mechanism in which the Laplace-dual conjugation $\diag(1/a)$
implements an exact duality between the \emph{time} of the spectral degeneracy and an
\emph{eigenvalue} of the dual pencil, collapsing a deficient Krylov space onto a double eigenvalue.
The result joins two previously separate facts about integrable MLZ --- that integrability forces an
exact crossing [Owusu--Wagh--Yuzbashyan, J.\ Phys.\ A \textbf{42}, 035206 (2009)], a statement about
the \emph{spectrum}, and the structure of the transition \emph{dynamics} --- and shows that
integrability, while it does not render the amplitude elementary, rigidifies it: the entire
monodromy data of the (non-elementary) connection problem becomes algebraic in the couplings.
\end{abstract}

\maketitle

\section{Introduction and motivation}
The multistate Landau--Zener problem asks for the matrix $P_{x\to j}=|\Scat_{xj}|^2$ of nonadiabatic
transition probabilities of $i\,\dot\psi=(H_0+tA)\psi$, $A=\diag(a)$, swept from $t=-\infty$ to
$+\infty$. Despite half a century of work it remains a \emph{constellation} of isolated exact
solutions --- Demkov--Osherov, bow-tie, Tavis--Cummings --- with little connecting theory. Two robust
facts anchor the landscape. First, the \emph{extreme}-slope survivals obey the universal
Brundobler--Elser law $P=\prod_m e^{-2\pi s_{nm}^2|a_n-a_m|}$ [Brundobler \& Elser, J.\ Phys.\ A
\textbf{26}, 1211 (1993)], with $s_{ij}=\gamma_i\gamma_j/(\eps_i-\eps_j)$. Second, the
\emph{interior} content --- the middle-slope survival $\Pmid$ (the diagonal $|\Scat_{mm}|^2$ of the
\emph{slope}-middle channel $m=\mathrm{argsort}(a)[1]$; the $\eps$-ordering fixes the basis but is not
itself a channel label) and one off-diagonal --- is the genuinely
open object, and for the minimal fully-coupled integrable case (Type-1, $N{=}3$, Cauchy/Gaudin
couplings $(H_0)_{ij}=\gamma_i\gamma_j(a_i-a_j)/(\eps_i-\eps_j)$) it is provably \emph{not} an
elementary or finite-special-function quantity: it is the connection constant of a rank-3
($GL_3$/$\mathfrak{sl}_3$) confluent isomonodromy problem --- the $W_3$ arena at self-dual central
charge $c=N{-}1=2$, not the rank-2 $c{=}1$ free-fermion (Painlev\'e) point [companion working paper,
this project].

Why is the interior hard, in one word? \emph{Accessory parameters.} Laplace-transforming the MLZ
equation turns $\Scat$ into the connection (Stokes) data of a linear ODE with a single irregular point
at infinity and a finite number of regular singular points. For Type-1 $N{=}3$ there is exactly one
finite singular point, and it is \emph{apparent}: the local solutions are single-valued
(indicial exponents $\{0,1,3\}$, no logarithm), so it carries no local physics --- only its
\emph{location} $v_*$ matters. That location is the accessory parameter. The accessory parameter is
the precise marker of the jump from solvable to transcendental: the hypergeometric equation has none
and its connection coefficients are closed-form (this is the ordinary two-level Landau--Zener answer,
built from Euler $\Gamma$ functions); adding one apparent singularity gives the Heun class, whose
connection coefficients have \emph{no} closed form, and the reason is exactly that the accessory
parameter is a free modulus on which the global connection data depend transcendentally. In a generic
accessory problem one does not even know \emph{where} the apparent singularity sits as a function of
the physical couplings; it is determined only by solving the whole problem.

This Letter settles the status of that one knob for integrable MLZ. We prove that for Type-1 $N{=}3$
the accessory parameter is \emph{pinned}: it equals, as an explicit rational function of
$(\gamma,\eps,a)$, the energy $E_*$ of the model's exact level crossing. The value of this is sharp.
It does not make $\Pmid$ elementary --- the no-go stands --- but it removes the single transcendental
degree of freedom that makes accessory problems hopeless, so that the connection problem one is left
with has \emph{completely explicit, algebraically-determined data}. In the language of the modern
Fredholm-determinant / Widom representation of isomonodromic connection constants [Cafasso, Gavrylenko \&
Lisovyy, arXiv:1712.08546], every input to the Riemann--Hilbert problem --- the formal monodromy
exponents and the accessory location alike --- is now a rational function of the couplings. Before
this result one could not even \emph{specify} the problem; after it, the only remaining unknown is the
genuinely transcendental connection constant itself.

\section{Setup: two objects on two sides}
Fix Type-1 $N{=}3$ data: $\eps_0<\eps_1<\eps_2$ real, $\gamma_i\neq0$ real, slopes $a_i$ distinct,
real, nonzero. Write $H(u)=H_0+uA$.

\emph{Object 1 --- the node (spectral side).} Let $\chi(u,E):=\det(EI-H(u))$, a cubic in both $u$ and
$E$. An exact level crossing is a point $(u_*,E_*)$ with $\chi=\partial_E\chi=0$: a double eigenvalue
$E_*$ at time $u_*$. That such a real crossing \emph{exists at all}, in apparent violation of the von
Neumann--Wigner non-crossing rule, is the theorem of Owusu, Wagh \& Yuzbashyan [J.\ Phys.\ A
\textbf{42}, 035206 (2009); arXiv:0807.0259]: any $H_0+uA$ possessing a nontrivial commuting partner
has such crossings, and Type-1 possesses one. For $N{=}3$ the crossing is unique, rational, and a
genuine \emph{node} of the spectral curve ($\partial_u\chi=0$ there as well: trivial monodromy).

\emph{Object 2 --- the accessory point (dynamical side).} The Laplace dual
$\psi_j(u)=\int_C e^{-iuv}B_j(v)\,dv$ obeys $B'(v)=-iM(v)B$ with $M(v):=\diag(1/a)(H_0-vI)$. With the
cyclic vector $e_0=(1,0,0)^{\!\top}$ define the cyclicity (Wronskian) determinant
\begin{equation}
W(v):=\det\!\big[\,e_0\ \big|\ M(v)e_0\ \big|\ M(v)^2 e_0\,\big].
\end{equation}
Because the first column is the $v$-independent $e_0$, $W(v)$ is \emph{linear} in $v$, so it has a
single root
\begin{equation}
v_*=-W_0/W_1\ \in\ \Z(\gamma^2,\eps,a),
\end{equation}
manifestly algebraic. This $v_*$ is the lone finite apparent singularity of the scalar third-order
connection ODE --- the accessory parameter.

The two objects live on opposite sides of the problem. $E_*$ is a feature of the \emph{instantaneous
spectrum} of $H(u)$ in the time domain --- a statement about eigenvalues. $v_*$ is a feature of the
\emph{singularity structure} of the transformed scattering problem in the Borel/Laplace plane --- a
statement about where an auxiliary ODE degenerates. There is no a priori reason they should be related.

\section{The theorem}
\begin{quote}
\textbf{Theorem (accessory $=$ node).} For every generic Type-1 $N{=}3$ model,
\begin{equation}
\boxed{\,v_*\ \equiv\ E_*\,}\qquad\text{as an identity in }\Z(\gamma,\eps,a).
\end{equation}
The lone finite apparent singularity of the connection ODE sits exactly at the spectral degeneracy.
\emph{Corollary:} the accessory parameter is the explicit rational function $E_*(\gamma,\eps,a)$, not
a free modulus; the connection problem has algebraically-fixed accessory data.
\end{quote}
(Genericity: $a_i$ distinct and nonzero, $\eps_i$ distinct, $\gamma_i\neq0$, and the node nondegenerate
so $W_1\neq0$.) Canonical instance ($\eps=(-2,0,3)$, $\gamma=(1,\tfrac45,\tfrac65)$,
$a=(-1,\tfrac12,2)$): $u_*=-187/750$, $v_*=E_*=-748/375$.

\section{Proof}
\emph{Route A (elimination identity).} $W(v)$ is linear (verified $\deg_v W=1$), giving
$v_*=-W_0/W_1$. The energies of a spectral degeneracy are the roots of
$\Phi(E):=\Res_u(\chi,\partial_E\chi)$. Substituting $E=v_*$ and clearing denominators, the numerator
reduces to the zero polynomial,
\begin{equation}
\Phi(v_*)\equiv0,
\end{equation}
equivalently $\Res_E(\Phi(E),\,W_1E+W_0)\equiv0$: $v_*$ \emph{is} a crossing energy. That it is the
\emph{node} branch (rather than a complex avoided-crossing branch) follows from Route B, which exhibits
the matching rational real $u_*$ with $\partial_E\chi=\partial_u\chi=0$. The identity is established
symbolically on the gauge slice $a=(s,1,2)$ with $\eps,\gamma$ fully symbolic and lifted to all generic
parameters by covariance (below); it is independently confirmed by exact rational arithmetic on
generic random samples.

\emph{Route B (structural mechanism --- the heart).} $W(v_*)=0$ says the $M(v_*)$-Krylov space of
$e_0$ is deficient: $e_0$ lies in a proper $M(v_*)$-invariant subspace. The link to the spectrum is a
one-line pencil identity:
\begin{equation}
\begin{aligned}
M(v)\,w=\mu\,w
  &\ \Longleftrightarrow\ (H_0-vI)\,w=\mu A\,w\\
  &\ \Longleftrightarrow\ H(-\mu)\,w=v\,w .
\end{aligned}
\label{eq:duality}
\end{equation}
An eigenpair $(\mu,w)$ of the dual matrix $M(v)$ is therefore an eigenpair of the \emph{physical}
Hamiltonian $H(u)$ at time $u=-\mu$ and energy $v$. The dual spectrum at fixed $v$ is the set of times
at which $v$ is an instantaneous eigenvalue. Hence the Krylov deficiency of $e_0$ at $v_*$ forces
$\chi(u,v_*)$, read as a polynomial in $u$, to have a \emph{double root} $u_*$
($\Disc_u\chi(\cdot,v_*)\equiv0$); at $(u_*,v_*)$ one finds $\chi=\partial_E\chi=0$ (double eigenvalue)
and $\partial_u\chi=0$ (genuine node). The rank-drop of the dynamical Wronskian and the spectral
degeneracy are \emph{the same event}, so $v_*=E_*$, with $u_*$ the rational node time delivered
explicitly. $\square$

\emph{Gauge covariance.} The construction is covariant under the affine slope group
$a_i\mapsto\alpha a_i+\beta$: $H_0$ depends only on differences $a_i-a_j$, so the group merely
relabels/rescales the MLZ time, and \emph{both} $v_*$ and $E_*$ transform identically (both are
energies on the same spectral curve). Thus an identity proven on a slice transverse to the orbits holds
everywhere. The Laplace dual forbids $a_i\to0$ (it uses $\diag(1/a)$), so the slice pins two nonzero
slopes ($a_1{=}1,a_2{=}2$) and keeps $a_0,\eps,\gamma$ general --- transverse to the $2$-parameter group
and reaching every generic gauge class.

\section{Why the spectral--dynamical link is nontrivial}
It is worth dwelling on why Eq.~(3) is not bookkeeping. The accessory parameter and the spectral
crossing are objects of fundamentally different type, attached to different stages of the problem and
living in different planes.

The crossing $E_*$ is a property of the \emph{time-domain generator}: it is where two instantaneous
eigenvalues of $H(u)$ touch as the control parameter $u$ is swept. It knows nothing, a priori, about
boundary conditions, scattering, or the analytic continuation that produces transition amplitudes ---
it is defined purely by the characteristic polynomial $\chi$.

The accessory point $v_*$ is a property of the \emph{frequency-domain solution operator}: it is where
the Laplace-transformed wave equation acquires an apparent singularity, a feature of the global
analytic structure of the scattering problem in the Borel plane. In the general theory of
isomonodromy and connection problems --- Heun, Painlev\'e, confluent Garnier --- the position of an
apparent singularity is the very definition of a modulus that is \emph{decoupled} from the local
spectral data: it is the free coordinate on the space of monodromy-preserving deformations, and its
dependence on the physical parameters is, generically, transcendental and implicit. One expects no
relation to any spectral invariant of a single operator; indeed the whole difficulty of Heun
connection theory is that this knob floats free.

That for Type-1 it does not float free, but is welded to the spectral degeneracy by an exact
algebraic identity, is a structural \emph{coincidence engineered by integrability}. The mechanism
\eqref{eq:duality} makes the coincidence visible: the Cauchy/Gaudin form of $H_0$, together with the
$\diag(1/a)$ conjugation intrinsic to the Laplace dual, sets up an exact duality that exchanges the
roles of \emph{time} and \emph{spectral parameter}. Under this duality the single linear Wronskian
$W(v)$ --- an object built from the dynamical side, the cyclic vector and the dual pencil --- encodes
the \emph{entire} time-domain spectral curve as seen from $e_0$, and its unique root is forced onto the
one point where that curve is doubly tangent: the node. The apparent singularity is thus the
\emph{shadow of the spectral degeneracy} cast through the Laplace transform. Strip away the Cauchy
structure (take a generic linear-in-time $N{=}3$ Hamiltonian) and the duality \eqref{eq:duality} and the
linearity of $W$ both break; the accessory parameter resumes floating, and no such identity can hold.

The bridge also has a conceptual payoff for the theory of integrable time-dependent systems. The
Owusu--Wagh--Yuzbashyan theorem is a statement on the spectral side: integrability $\Rightarrow$ an
exact crossing. Our theorem propagates integrability to the dynamical side: that same crossing
$\Rightarrow$ a pinned accessory parameter $\Rightarrow$ algebraic monodromy data. Chaining them,
\begin{equation*}
\begin{aligned}
\text{commuting partner}
  &\ \Rightarrow\ \text{exact crossing}\\
  &\ \Rightarrow\ \text{accessory parameter pinned}\\
  &\ \Rightarrow\ \text{algebraic monodromy data},
\end{aligned}
\end{equation*}
gives a precise answer to a question the constellation of solvable models leaves implicit: \emph{what
does integrability buy when it does not buy solvability?} It buys rigidity. The transition amplitude
remains transcendental --- integrability does not lower it to elementary functions (the ``Abelian
ceiling'') --- but its analytic skeleton, the monodromy data of the governing connection problem, is
forced to be an algebraic function of the couplings. ``Integrable but not solvable'' thereby acquires
a sharp meaning: these are exactly the models whose amplitudes are transcendental \emph{with
algebraically-determined data}.

\section{Avenues opened and unblocked}
The theorem is enabling rather than terminal; several lines that were previously ill-posed become
concrete.

\emph{(i) The connection constant is now fully specified.} With the formal monodromy exponents
$\Theta=\diag(c_i)$, $c_i=\sum_{j\ne i}s_{ij}^2(a_i-a_j)$ (established independently by the Wasow
recursion in the Laplace frame) and the accessory location $v_*=E_*$ both algebraic, every datum of
the rank-3 Riemann--Hilbert / Fredholm-determinant problem for $\Pmid$ is an explicit rational
function of $(\gamma,\eps,a)$. The remaining unknown is the genuinely transcendental connection
constant alone; computing or bounding it (e.g.\ via the Widom/Fredholm representation, or a Nekrasov-type
expansion) is now a well-posed problem with no free parameters.

\emph{(ii) Constructive exact-WKB with explicit turning-point data.} The accessory point is the
location of the spectral-network joint that carries the irreducible content of $\Pmid$. Knowing it in
closed form pins the joint and the local Weber/parabolic-cylinder connection pieces, unblocking the
constructive ``two-shear $\times$ junction'' product as an object with explicit ingredients rather than
implicit ones.

\emph{(iii) An algebraic equation for the elementary boundary.} The model degenerates to elementary
(a level decouples) precisely where the accessory point collides with the irregular point at infinity,
i.e.\ on the locus $W_1=0$. Because $v_*=E_*$ is algebraic, this boundary --- the edge of solvability
in parameter space --- is now cut out by an explicit polynomial condition, sharpening the classifier
``elementary $\Leftrightarrow$ a level decouples'' into an equation.

\emph{(iv) A spectral duality to generalize.} The pencil identity \eqref{eq:duality} is not special to
$N{=}3$: for general $N$ it still exchanges time and dual eigenvalue. This suggests a conjecture --- for
the integrable Type-1 family at all $N$, the accessory variety of the connection problem coincides with
the spectral-degeneracy variety of $H(u)$ --- and a strategy (Krylov rank-drop $\Leftrightarrow$
spectral tangency) to attempt it. If true, it would extend ``algebraic monodromy data'' to the whole
integrable Type-1 hierarchy and turn a single-model result into a structural theorem about the boundary
of the solvable constellation.

\emph{(v) A theorem-backed physical handle.} The node is a real-time exact degeneracy whose location
$u_*$ is now an explicit rational function of the couplings, and which provably controls the analytic
structure of the middle transition. This makes it a concrete target for studying --- or engineering ---
the dynamics of the interior channel: the one event that organizes the hard part of the amplitude is no
longer hidden but addressable.

\emph{(vi) Owed finish.} The strict no-logarithm certificate for the apparent point (the vanishing of
the order-$2$ Frobenius coefficient as a single fully-symbolic identity) is the one secondary item not
yet closed; the indicial signature $\{0,1,3\}$ is proven. This is a finishing computation, not an
obstruction.

\begin{acknowledgments}
The exact level crossing underpinning this identity is the theorem of Owusu, Wagh \& Yuzbashyan
(2009); this Letter establishes its dynamical counterpart. Proof artifacts and an exhaustive
verification record accompany the companion working paper.
\end{acknowledgments}

\end{document}
